\newpage
\recipe{Risotto Patrio}{Perfezionato, nella mia non modesta opinione, dalla \href[]{https://www.gadda.ed.ac.uk/Pages/resources/essays/risotto.php}{versione di Carlo Emilio Gadda}. Per 4 commensali}{}{}

\begin{multicols}{2}
	\subsection*{Ingredienti}
	\begin{ingredients}
		\item 1 cipolla tritata
		\item 350 g di riso per risotto (e.g., di tipo Vialone)
		\item Burro Lodigiano di classe, \textit{quantum prodest}
		\item 1.5 L di brodo di carne saporito
		\item 0.2 g (o 50 stimmi) di zafferano
		\item Vino bianco per sfumare
		\item Parmigiano Reggiano grattuggiato
		\item Funghi porcini, chanterelles, o spugnole
	\end{ingredients}
	
	\columnbreak
	
	\subsection*{Utensili}
	\begin{equipment}
		\item Casseruola di rame stagnato
		\item Ramaiolo
		\item Mestolo di legno
	
	\end{equipment}
	\vfill
	\end{multicols}

\medskip

% --- Preparation ---
In una scodella marginale, preparare 1.5 L di \textbf{brodo} di carne \textbf{zafferanato}, da avere in pronto. 
Se disponibili, aggiungere al brodo anche i \textbf{funghi}. 
Nella casseruola, sciogliere del \textbf{burro}\footnote{E, volendo, un cucchiaio colmo di olio d'oliva perchè il burro bruci meno facilmente.} e la \textbf{cipolla} in minimi pezzi.
Al primo soffriggere del connubio butirroso-cipollino buttare il \textbf{riso}, da tostare e rosolare contro il fondo stagnato. 
Sfumare col \textbf{vino} bianco.
Iniziare quindi la bollitura del riso, versando in piccole dosi il brodo fino a cottura avvenuta: che il riso non sia scotto, ma poco piu che al dente. 
Tolto il risotto dal fuoco, \textbf{mantecarlo} e aggiungere tanto \textbf{parmigiano} quanto piace. 
Coprire con un coperchio e lasciar riposare per 10 minuti prima di servirlo con del pane rustico e un bicchiere di vino rosso. 


\vspace*{0.75 cm}

% --- Description ---
\begin{recipeblock}
	
\end{recipeblock}